% !PRINTED 203
\weberpage{203}

%<<CONT
holds, then in it the variable $v$ may be replaced by $v + 4\mu K + 4\mu' iK'$.
On account of the periodicity of sn\,$v$, cn\,$v$, dn\,$v$ it follows from this,
however, that the equation

\[
F(x,\ \mathrm{sn}\,v,\ \mathrm{cn}\,v,\ \mathrm{dn}\,v) = 0
\]

is satisfied for \gesperrt{all} roots of equation (4). In order to determine its \gesperrt{Galois}
group, we first fix the \gesperrt{domain of rationality}. It is to comprise the
following quantities:

\begin{enumerate}
\item[1.] rational numbers,
\item[2.] rational functions of $\varkappa^2$,
\item[3.] the three functions sn\,$v$, cn\,$v$, dn\,$v$,
\item[4.] the quantities $\mathrm{sn}\left(\dfrac{4\mu K' + 4\mu' iK'}{n}\right)$. % [sic]
\end{enumerate}

From (2) it then follows that

\[
\mathrm{cn}\left(\frac{4\mu K + 4\mu' iK'}{n}\right), \qquad \mathrm{dn}\left(\frac{4\mu K + 4\mu' iK'}{n}\right)
\]

also belong to the domain of rationality.

We now remark that, by the addition theorem, each of the quantities
$x_{\mu,\mu'}$ is rationally expressible through every other one among
them. If in particular

\weq{5}{x_{\mu,\mu'} = F_{\mu,\mu'}(x_{0,0})}

holds, then

\weq{6}{x_{\mu+\nu,\mu'+\nu'} = F_{\mu,\mu'}(x_{\nu,\nu'}) = F_{\mu+\nu,\mu'+\nu'}(x_{0,0}).}

From this it follows that the \gesperrt{Galois} group of our equation consists of the
$n^2$ permutations $S_{\nu,\nu'}$ which one obtains by increasing, in the
roots $x_{\mu,\mu'}$, the indices $\mu, \mu'$ all by the same pair of
numbers $\nu, \nu'$ (where each index is to be taken according to the
modulus $n$).

For by (5) every rational function of the roots
$x_{\mu,\mu'}$ can be expressed rationally through $x_{0,0}$ in the form

\[
\Phi(x_{0,0})
\]

and therefore passes over, by (6), under the substitution
$S_{\nu,\nu'}$ into $\Phi(x_{\nu,\nu'})$; if it thus remains unchanged under this
substitution, then \gesperrt{it} is \gesperrt{rational}. Conversely it follows from the
irreducibility that every rational equation between the roots,
since it can be put into the form:

\[
\psi(x_{0,0}) = 0,
\]

and thus

\[
\psi(x_{\nu,\nu'}) = 0
\]
%CONT>>
